% =============================================================================
% 文件: sections/s2_analysis.tex   第二章 问题分析
% =============================================================================
\section{问题分析}

\subsection{总体思路}

本题的物理内核是一个\textbf{圆柱体内的耦合传热传质问题}：温度场决定水分扩散系数，
含水率场又通过密度、比热、导热系数反过来影响温度场，二者通过非线性物性双向耦合。
四个问题共用同一套控制方程与同一个求解器，区别只在于物性关系、几何是否随时间变化
以及输出要求。因此本文采取“\textbf{一个模型 + 一个求解器 + 四组配置}”的路线：
先把控制方程与离散格式写死，再用数据检验边界条件与平台假设，最后按题面要求输出
四套结果。这样做的直接好处是：任何一处改进（例如边界重构、时间步收敛）都会自动
惠及全部四问，也便于定位问题。

总体技术路线为：

\begin{center}
\fbox{\parbox{0.94\textwidth}{
数据探索（附件 1 分段与平台、附件 2 收缩轨迹、量级检查）
$\rightarrow$ 建立一维轴对称耦合传热传质方程组
$\rightarrow$ 节点中心有限体积离散 + 半单元表面重构
$\rightarrow$ 隐式 Euler + Picard 迭代求解
$\rightarrow$ 解析解/守恒律/多方法互证
$\rightarrow$ 灵敏度与不确定度量化
$\rightarrow$ 四问结果输出与论文撰写}}
\end{center}

\subsection{为什么选择“一维轴对称 + 有限体积 + 隐式时间推进”}

\paragraph{机理层面：为什么可以降为一维径向}
药材长 $25\ \mathrm{cm}$、半径 $2\ \mathrm{cm}$，长径比 $12.5:1$，且热风沿轴向吹过时
沿长度的温度与湿度梯度远小于径向梯度。端面虽然会额外失水，但其面积占比为
$2\pi R_0^2/(2\pi R_0 L)=R_0/L=8\%$，且端面附近同样是“表面先干、内部后干”的
一维径向结构。因此忽略轴向与端面效应的系统误差量级约为 $8\%$ 的表面积偏差，
这一量级小于本题传质系数与空气含湿量口径所引入的不确定度（见第七章），
属于可接受的简化。

\paragraph{模型形式：为什么是抛物型扩散方程而不是经验薄层模型}
常见的薄层干燥模型（如 Page、Midilli 模型）把整颗物料当作“一个点”，
只能给出平均含水率随时间的变化，无法回答题面要求的“到药材中心不同距离处的
含水率”，也无法给出“各处都低于 $0.15$”的判定。Fick 第二定律形式的扩散方程
则天然给出空间分布，并且其圆柱几何的解析解（Crank 解，Bessel 级数）可作为
数值实现的独立校验基准。因此本文选择分布参数模型而非集总参数模型。

\paragraph{耦合方式：为什么温度场必须与含水率场联立}
附录 3、附录 4 中 $D=2.4\times10^{-3}\mathrm{e}^{-0.45/C}\mathrm{e}^{-3850/T}$
同时依赖 $C$ 与 $T$。若把 $T$ 固定为环境温度，则在 $50\ ^\circ\mathrm{C}$ 附近
$\mathrm{e}^{-3850/T}$ 的取值误差会直接放大到 $D$ 上；而 $T$ 的实际演化（表面
先升、中心后升）在模型里恰好决定了“哪一段时间扩散得快”。因此必须联立求解。
反之，含水率通过 $\rho(C)$、$c_p(C)$、$k(C)$ 影响热容与导热，也不能单方面固定。

\paragraph{离散方法：为什么是有限体积而不是有限差分}
本题的关键困难是\textbf{扩散系数在空间上跨越数量级}：由
$D\propto \mathrm{e}^{-0.45/C}$，在 $C$ 从 $2.55$ 降到 $0.15$ 的过程中
$D$ 下降约一个数量级，且这个陡降集中发生在表面附近很薄的一层内。
有限体积法在界面上用\textbf{调和平均}扩散系数，等价于把两个相邻控制体视为
串联热阻/质阻，可以严格保证界面通量的守恒性与非负性；而直接对
$\partial/\partial r(rD\partial C/\partial r)$ 做中心差分会在大 $D$ 梯度处
产生非物理的“伪通量”。这一点在第六章的守恒检验中得到证实
（水量收支闭合误差为 $0$）。

\paragraph{时间推进：为什么是隐式 Euler + Picard 迭代}
物性对解的依赖是强非线性的（指数型），显式格式的时间步受
$\Delta t\lesssim \Delta r^2/(2D)$ 的稳定性限制，在 $N=1600$、$D\sim10^{-8}$ 时
要求 $\Delta t<10^{-3}\ \mathrm{s}$，对于 $57\ \mathrm{h}$ 的过程不可接受。
隐式 Euler 无条件稳定，代价是每一步需要求解非线性代数方程组；本文用
Picard（不动点）迭代线性化：用上一次迭代的温度与含水率计算全部物性，
组装三对角方程组后求解。这样每个时间步的机时只与 $N$ 成正比，
$57\ \mathrm{h}$ 的全过程算例在普通笔记本上只需几分钟。

\paragraph{问题 4 的坐标选择：为什么用物质坐标}
收缩域的难点是边界在动。若在物理坐标 $r$ 下直接追踪移动边界，网格必须随时间
移动，方程中要额外保留“网格速度乘梯度”的对流项，稍有疏忽就会引入虚假的
对流扩散。本文改用物质坐标 $\xi=r/R(t)$：把变动域 $[0,R(t)]$ 映射到固定区间
$[0,1]$，并且让 $\xi$ 恰好标记\textbf{径向材料点}。在纯径向仿射收缩假设下，每个材料点
的 $\xi$ 恒为常数，于是时间导数 $\partial/\partial t|_\xi$ 就是物质导数，
原始的守恒形式 $\partial C/\partial t=\nabla\cdot(D\nabla C)$ 在
$\xi$ 坐标下变为
$\partial C/\partial t|_\xi=(1/\xi R^2)\partial_\xi(\xi D\partial_\xi C)$，
\textbf{不需要附加对流项}。这一步是问题 4 建模正确性的关键。

\subsection{数据探索：从数据特征反推模型结构}

建模前的数据探索不是“看一眼数据”，而是要回答三个会改变模型结构的问题。

\paragraph{（1）附件 1 的环境是“一个阶段”还是“两个阶段”？平台值取多少？}
附件 1 的 241 个点（$60\ \mathrm{s}$ 间隔，$0\sim14400\ \mathrm{s}$）中，
温度由 $28.000\ ^\circ\mathrm{C}$ 升至约 $50\ ^\circ\mathrm{C}$，水分浓度由
$0.01963$ 升至约 $0.0500$。用 $600\ \mathrm{s}$ 滑窗线性斜率作为判据
（$|\mathrm{d}T/\mathrm{d}t|<0.0017\ ^\circ\mathrm{C/s}$），可得预热段约为
$0\sim3420\ \mathrm{s}$，其后进入平台段，末段 $3600\ \mathrm{s}$ 的温度标准差
仅 $0.1535\ ^\circ\mathrm{C}$，说明平台确实存在而非缓慢漂移。
但题面的干燥过程长达 $2\sim3$ 天，附件 1 只覆盖前 $4\ \mathrm{h}$，
因此 $t>14400\ \mathrm{s}$ 的环境\textbf{必须由模型假设给出}。
本文比较了两种口径：
（i）实测口径 —— 末段 $3600\ \mathrm{s}$ 的算术平均
$\Tair=49.9957\ ^\circ\mathrm{C}$、$\Cair=0.049991$；
（ii）外推口径 —— 用一阶惯性模型
$x(t)=x_\infty-(x_\infty-x_0)\mathrm{e}^{-t/\tau}$ 拟合并取其渐近值
$\Tair=50.1533\ ^\circ\mathrm{C}$、$\Cair=0.050653$（以实际拟合输出为准）。
必须指出：口径（ii）的拟合值仍是\textbf{外推而非实测}，
属于外推而非常测。按“数据优先”的原则，本文以口径（i）作为基准，
把口径（ii）以及其他窗口长度放进灵敏度分析，量化平台口径对 $t_3$、$t_4$ 的影响。

\paragraph{（2）附件 2 的半径历史能否与密度经验式自洽？}
附件 2 的半径由 $2.000\ \mathrm{cm}$ 单调降至 $1.198\ \mathrm{cm}$，
二阶差分最大值达 $0.048\ \mathrm{cm}$，说明轨迹存在折点，插值方式会影响结果。
更重要的是：如果药材的收缩完全由失水引起且干物质不可压缩，则单位长度干物质质量
$m_s=\rho_{\mathrm{d}}\pi R^2$（其中干物质表观密度
$\rho_{\mathrm{d}}=\rho(C)/(1+C)$）应守恒。用附录 4 的 $\rho(C)$ 与附件 2 的
$R(t)$ 分别计算，$m_s$ 的变化量达百分之十几（详见 7.6 节）。
这说明\textbf{附件 2 的实测收缩轨迹与附录 4 的密度经验式不是由同一个固体守恒方程
标定的}。因此问题 4 的正确定位是“\textbf{给定实测收缩轨迹下的预测}”，
而不是完整的力学收缩模型；本文在论文中明确声明这一限制，并在检验章节量化它。

\paragraph{（3）量纲与数量级是否合理？}
附录 3、附录 4 的 $\mathrm{e}^{-3850/T}$ 中的 $T$ 必须用\textbf{开尔文}：
$T=301\ \mathrm{K}$ 时该因子为 $2.79\times10^{-6}$，$T=323\ \mathrm{K}$ 时为
$6.66\times10^{-6}$；若误用摄氏度则指数变成 $\mathrm{e}^{-3850/50}=10^{-34}$，
扩散系数将小到物理上不可能。按 SI 单位计算，附录 3 的 $D$ 在
$C=2.55$、$T=50\ ^\circ\mathrm{C}$ 时为 $1.35\times10^{-8}\ \mathrm{m^2/s}$，
在 $C=0.15$ 时降至 $3.35\times10^{-10}\ \mathrm{m^2/s}$，落在农产食品热风干燥有效水分
扩散系数的文献区间（约 $10^{-11}\sim10^{-8}\ \mathrm{m^2/s}$）内，量级正确。
按 $R^2/D$ 估计的特征扩散时间约为 $8\sim332\ \mathrm{h}$（取
$C\in[0.15,2.55]$、$T\in[28,51]\ ^\circ\mathrm{C}$ 的端点）。这是整个物性取值
盒上的保守范围；在干燥主体区间内的代表值约为十几至数十小时，与题面
“干燥持续 $2\sim3$ 天”在数量级上一致。

\begin{figure}[H]
\centering
\includegraphics[width=0.88\textwidth]{fig01_room_environment.png}
\caption{附件 1 烘房环境与工艺分段：上图为温度、下图为水分浓度，
点线为滑窗斜率判据给出的预热/恒温分界（约 3420 s），
虚线为末段 3600 s 的实测均值（本文基准平台取值）}
\label{fig:room}
\end{figure}

\begin{figure}[H]
\centering
\includegraphics[width=0.88\textwidth]{fig02_radius_history.png}
\caption{附件 2 药材半径历史：(a) 半径轨迹；(b) 相对体积收缩 $V/V_0$；
(c) 收缩速率逐段递减，说明收缩在干燥后期趋于停止}
\label{fig:radius}
\end{figure}

\begin{figure}[H]
\centering
\includegraphics[width=0.88\textwidth]{fig05_plateau_estimators.png}
\caption{平台边界的两种估计口径：(a) 温度；(b) 水分浓度。
虚线为一阶惯性拟合曲线及其渐近值，实线为末段实测均值。
注意拟合渐近的水分浓度已超出附件 1 的实测最大值，属于外推口径}
\label{fig:plateau}
\end{figure}

\subsection{各问的求解路线}

\begin{itemize}[leftmargin=2.2em,itemsep=0.2em]
  \item \textbf{问题 1}：$R$ 固定、物性取附录 2，环境边界对附件 1 做分段线性插值。
        由于 $0\sim1800\ \mathrm{s}$ 完全落在附件 1 的覆盖区间内，本问\textbf{不涉及
        任何平台假设}，是模型最“干净”的验证场景。
  \item \textbf{问题 2}：$R$ 固定、物性取附录 3，环境在 $t\le14400\ \mathrm{s}$ 用
        附件 1 插值、$t>14400\ \mathrm{s}$ 用平台假设；输出 $3\ \mathrm{h}$ 内的网格结果。
  \item \textbf{问题 3}：与问题 2 同一模型，推进到 $\max_r C<0.15$；
        干燥时间 $t_3$ 由最大含水率曲线与 $0.15$ 的线性插值交点确定，
        并给出 $\Delta t$ 与 $N$ 双重收敛证据。
  \item \textbf{问题 4}：物性取附录 4，几何用物质坐标，半径由附件 2 插值；
        同一判据给出 $t_4$，并对比线性插值与保形三次插值（PCHIP）的差异。
\end{itemize}
